\documentclass[a4paper, final]{article}

% ============================================================
% НАСТРОЙКИ
% ============================================================
\usepackage{fontspec}
\setmainfont{Times New Roman}
\setmonofont{Courier New}
\usepackage[14pt]{extsizes}

\usepackage{polyglossia}
\setdefaultlanguage{russian}

\usepackage[left=20mm, top=20mm, right=20mm, bottom=20mm, footskip=10mm]{geometry}
\usepackage{ragged2e}
\usepackage{indentfirst}
\usepackage{hyperref}
\usepackage{graphicx}
\usepackage{tikz}
\usepackage{caption}
\usepackage{subcaption}
\usepackage{xcolor}
\usepackage{listings}
\usepackage{algorithm}
\usepackage{algorithmic}
\usepackage{soul}
\usepackage{float}

\definecolor{codegray}{gray}{0.95}
\definecolor{codered}{rgb}{0.6,0,0}
\definecolor{codeblue}{rgb}{0,0,0.6}
\definecolor{codegreen}{rgb}{0,0.5,0}

\lstdefinestyle{common}{
  basicstyle=\ttfamily\small,
  backgroundcolor=\color{codegray},
  frame=single,
  numbers=left,
  numberstyle=\tiny\color{gray},
  stepnumber=1,
  numbersep=6pt,
  showspaces=false,
  showstringspaces=false,
  showtabs=false,
  tabsize=2,
  breaklines=true,
  breakatwhitespace=true,
  captionpos=b,
  keywordstyle=\color{codeblue}\bfseries,
  commentstyle=\color{codegreen}\itshape,
  stringstyle=\color{codered}
}

\lstdefinestyle{cstyle}{
    style=common,
    language=C++,
    inputencoding=utf8,
    extendedchars=true,
    keepspaces=true,
    columns=fullflexible,
    literate={а}{{\selectfont\char224}}1
             {б}{{\selectfont\char225}}1
             {в}{{\selectfont\char226}}1
             {г}{{\selectfont\char227}}1
             {д}{{\selectfont\char228}}1
             {е}{{\selectfont\char229}}1
             {ё}{{\"e}}1
             {ж}{{\selectfont\char230}}1
             {з}{{\selectfont\char231}}1
             {и}{{\selectfont\char232}}1
             {й}{{\selectfont\char233}}1
             {к}{{\selectfont\char234}}1
             {л}{{\selectfont\char235}}1
             {м}{{\selectfont\char236}}1
             {н}{{\selectfont\char237}}1
             {о}{{\selectfont\char238}}1
             {п}{{\selectfont\char239}}1
             {р}{{\selectfont\char240}}1
             {с}{{\selectfont\char241}}1
             {т}{{\selectfont\char242}}1
             {у}{{\selectfont\char243}}1
             {ф}{{\selectfont\char244}}1
             {х}{{\selectfont\char245}}1
             {ц}{{\selectfont\char246}}1
             {ч}{{\selectfont\char247}}1
             {ш}{{\selectfont\char248}}1
             {щ}{{\selectfont\char249}}1
             {ъ}{{\selectfont\char250}}1
             {ы}{{\selectfont\char251}}1
             {ь}{{\selectfont\char252}}1
             {э}{{\selectfont\char253}}1
             {ю}{{\selectfont\char254}}1
             {я}{{\selectfont\char255}}1
             {А}{{\selectfont\char192}}1
             {Б}{{\selectfont\char193}}1
             {В}{{\selectfont\char194}}1
             {Г}{{\selectfont\char195}}1
             {Д}{{\selectfont\char196}}1
             {Е}{{\selectfont\char197}}1
             {Ё}{{\"E}}1
             {Ж}{{\selectfont\char198}}1
             {З}{{\selectfont\char199}}1
             {И}{{\selectfont\char200}}1
             {Й}{{\selectfont\char201}}1
             {К}{{\selectfont\char202}}1
             {Л}{{\selectfont\char203}}1
             {М}{{\selectfont\char204}}1
             {Н}{{\selectfont\char205}}1
             {О}{{\selectfont\char206}}1
             {П}{{\selectfont\char207}}1
             {Р}{{\selectfont\char208}}1
             {С}{{\selectfont\char209}}1
             {Т}{{\selectfont\char210}}1
             {У}{{\selectfont\char211}}1
             {Ф}{{\selectfont\char212}}1
             {Х}{{\selectfont\char213}}1
             {Ц}{{\selectfont\char214}}1
             {Ч}{{\selectfont\char215}}1
             {Ш}{{\selectfont\char216}}1
             {Щ}{{\selectfont\char217}}1
             {Ъ}{{\selectfont\char218}}1
             {Ы}{{\selectfont\char219}}1
             {Ь}{{\selectfont\char220}}1
             {Э}{{\selectfont\char221}}1
             {Ю}{{\selectfont\char222}}1
             {Я}{{\selectfont\char223}}1,
}

\linespread{0.98}

\definecolor{link_color}{HTML}{0000EE}

\hypersetup{
    colorlinks=true,
    linkcolor=link_color,
    urlcolor=link_color,
    citecolor=link_color
}

% ============================================================
% ДОКУМЕНТ
% ============================================================
\begin{document}

% ============================================================
% ТИТУЛЬНЫЙ ЛИСТ
% ============================================================
{
\thispagestyle{empty}
\begin{center}
    \hfill \break
    \hfill \break
    МИНИСТЕРСТВО НАУКИ И ВЫСШЕГО ОБРАЗОВАНИЯ РОССИЙСКОЙ ФЕДЕРАЦИИ\\[5 pt]
    Федеральное государственное автономное образовательное учреждение высшего образования
    «Санкт-Петербургский политехнический университет Петра Великого»\\[10 pt]
    Институт компьютерных наук и кибербезопасности\\ [5 pt]
    Высшая школа технологий искусственного интеллекта \\ [5 pt]
    Направление: 02.03.01 Математика и компьютерные науки\\[30pt]

    {\large ОТЧЕТ ПО Лабораторной работе №4}\\[10 pt]
    {\large по дисциплине Программирование микроконтроллеров}\\[40 pt]
    {\large\bf Использование таймера и аппаратных прерываний.}
\end{center}

\vspace{110 pt}

\begin{tabular*}{460pt}{@{\extracolsep{\fill}} l r}
     Студент,\tabularnewline группы 5130201/40003 &  \line(1,0){100} \quad Четвергов И. C.\\
     & \\
     Руководитель,\tabularnewline Преподаватель &  \line(1,0){100} \quad Вербова Н. М.\\
     & \\
     & \\
     & \\
     & \guillemotleft \line(1,0){30} \guillemotright \line(1,0){100} \, 2026 \,г.

\end{tabular*} \\

\vfill

\begin{center}
    Санкт-Петербург, 2026
\end{center}
\newpage
}

\tableofcontents
\newpage

% ============================================================
% ОСНОВНОЕ СОДЕРЖАНИЕ
% ============================================================

\section{Тема}
Использование таймера для формирования заданного временного интервала и аппаратного прерывания для перевода микроконтроллера в режим пониженного энергопотребления.

\section{Цель работы}
Ознакомиться с основными методами формирования заданных интервалов времени и перевода микроконтроллера в режим пониженного энергопотребления. Закрепить навыки работы с осциллографом и оценочной платой MCBSTM32F200 в качестве измерительного генератора.

\section{Постановка задачи}
Используя библиотеки Keil~µVision5, разработать программу для микроконтроллера (МК) STM32F200, которая при помощи таймера формирует периодическое попеременное включение и выключение светодиодов PG6 и PG7 с заданными временными характеристиками. При нажатии на кнопку «WAKEUP» программа должна переводить МК в спящий режим, а при её отпускании~--- пробуждать МК.

\section{Теоретические данные}

Для формирования и измерения временных интервалов в МК применяется аппаратная подсистема~--- таймер. Таймер автоматически отсчитывает начало и конец заданного процесса или его отдельных этапов.

МК STM32F200 оснащён восемью периферийными программируемыми 16-разрядными таймерами TIM1–TIM8 с перезагружаемыми счётчиками и программируемыми предделителями. Они применяются для измерения длительности входных импульсов и генерации прямоугольных сигналов. С помощью предделителей и контроллера тактирования RCC длительность импульсов и период их следования можно регулировать от нескольких микросекунд до нескольких миллисекунд. Таймеры работают независимо друг от друга и не потребляют общих ресурсов, но при необходимости могут синхронизироваться.

Помимо периферийных таймеров, в ядре микроконтроллера имеется системный таймер SysTick, изначально предназначенный для формирования временной базы ОСРВ. Тем не менее его периодические прерывания допускается использовать и в других целях. В микроконтроллерах на ядре Cortex время реакции на прерывание строго фиксировано.

Для настройки SysTick в файле \texttt{core\_cm3.h} предусмотрена функция \texttt{SysTick\_Config(uint32\_t ticks)}, принимающая коэффициент деления тактовой частоты. По умолчанию тактовая частота ядра составляет 12~МГц.

При переходе в режим сна или глубокого сна тактирование процессора Cortex-M3 останавливается; в режиме глубокого сна дополнительно отключаются флэш-память и основной тактовый генератор. Управление этими режимами осуществляется через регистр \texttt{SCB\_SCR}, а конфигурация режима сна задаётся с помощью маски \texttt{SCB\_SCR\_SLEEPONEXIT\_Msk} из файла \texttt{core\_cm3.h}.

\section{Алгоритм программы}

\subsection{Инициализация периферии}
\begin{itemize}
    \item Включение тактирования портов GPIOG и GPIOA через регистр \texttt{RCC->AHB1ENR}.
    \item Включение тактирования контроллера конфигурации системы SYSCFG через регистр \texttt{RCC->APB2ENR}.
    \item Настройка выводов PG6, PG7 как цифровых выходов; PA0~--- как цифрового входа.
\end{itemize}

\subsection{Настройка внешних прерываний (EXTI)}
\begin{itemize}
    \item Разрешение прерываний по линии EXTI0 через регистр \texttt{EXTI->IMR}.
    \item Настройка срабатывания по нарастающему фронту (\texttt{EXTI->RTSR}).
    \item Привязка вывода PA0 к линии прерывания через \texttt{SYSCFG->EXTICR[0]}.
\end{itemize}

\subsection{Настройка системного таймера SysTick}
\begin{itemize}
    \item Вызов \texttt{SystemCoreClockUpdate()} для получения актуальной тактовой частоты.
    \item Конфигурация SysTick с периодом прерывания 1~мс через \texttt{SysTick\_Config}.
\end{itemize}

\subsection{Режим пониженного энергопотребления}
\begin{itemize}
    \item Установка бита \texttt{SLEEPDEEP} в \texttt{SCB->SCR} для выбора режима глубокого сна.
    \item Настройка регистров \texttt{PWR->CR} для выбора режима standby.
\end{itemize}

\subsection{Управление прерываниями (NVIC)}
\begin{itemize}
    \item Задание групп и приоритетов прерываний от EXTI0 и SysTick.
    \item Активация обработчиков прерываний.
\end{itemize}

\subsection{Основной цикл}
В бесконечном цикле поочерёдно вызываются функции управления светодиодами: включается PG6, затем PG7, после чего цикл повторяется.

\subsection{Обработчик прерывания EXTI0}
При нажатии кнопки «WAKEUP» выполняется программная задержка для устранения дребезга, затем инвертируется бит \texttt{SLEEPONEXIT}~--- МК переходит в спящий режим или пробуждается. Фронт срабатывания прерывания переключается на противоположный.

\newpage

\section{Код программы}

\begin{lstlisting}[style=cstyle, caption={MyBlinky.c}]
#include "stm32f207xx.h"
#include "stdbool.h"
#include "stdio.h"

// Программная задержка на основе пустого цикла
void delay(unsigned long cycles)
{
    unsigned long i = 0;
    for (i = 0; i < cycles; i++)
    { }
}

// Счётчик миллисекунд, обновляемый системным таймером
volatile uint32_t msTicks = 0;

// Задержка в миллисекундах на основе аппаратного таймера
void delay_ms(uint32_t ms)
{
    uint32_t t = msTicks + ms;
    while (t > msTicks);
}

// Обработчик прерывания системного таймера SysTick — увеличивает счётчик на 1
void SysTick_Handler()
{
    msTicks++;
}

// Обработчик внешнего прерывания по линии EXTI0 (кнопка WAKEUP)
void EXTI0_IRQHandler(void)
{
    // Программная задержка для подавления дребезга контактов
    delay(5000);

    // Переключение бита SLEEPONEXIT: переход в сон или пробуждение МК
    SCB->SCR ^= SCB_SCR_SLEEPONEXIT_Msk;

    // Инверсия фронта срабатывания прерывания
    EXTI->RTSR ^= EXTI_RTSR_TR0;
    EXTI->FTSR ^= EXTI_FTSR_TR0;

    // Сброс флага ожидающего прерывания
    EXTI->PR |= EXTI_PR_PR0;
}

// Попеременное включение светодиодов PG6 и PG7
void DiodPG6(void)
{
    // Включить светодиод на PG6
    GPIOG->ODR |= 1ul << 6;
    delay(500000);
    // Выключить светодиод на PG6
    GPIOG->ODR &= ~1ul << 6;

    // Включить светодиод на PG7
    GPIOG->ODR |= 1ul << 7;
    delay(500000);
    // Выключить светодиод на PG7
    GPIOG->ODR &= ~1ul << 7;
}

// Заготовка для дополнительного управления светодиодом PG7
void DiodPG7(void)
{

}

int main()
{
    // Включить тактирование порта GPIOG
    RCC->AHB1ENR |= 1ul << 6;
    // Включить тактирование порта GPIOA
    RCC->AHB1ENR |= 1ul << 0;
    // Включить тактирование контроллера конфигурации системы SYSCFG
    RCC->APB2ENR |= 1ul << 14;

    // Настроить PG6 как цифровой выход
    GPIOG->MODER = (GPIOG->MODER & ~(1ul << 13)) | 1ul << 12;
    // Настроить PG7 как цифровой выход
    GPIOG->MODER = (GPIOG->MODER & ~(1ul << 15)) | 1ul << 14;
    // Настроить PG8 как цифровой выход
    GPIOG->MODER = (GPIOG->MODER & ~(1ul << 17)) | 1ul << 16;
    // Настроить PA0 как цифровой вход (кнопка WAKEUP)
    GPIOG->MODER &= ~GPIO_MODER_MODER0;

    // Разрешить прерывания по внешней линии EXTI0
    EXTI->IMR |= EXTI_IMR_MR0_Msk;
    // Настроить срабатывание по нарастающему фронту (нажатие кнопки)
    EXTI->RTSR |= EXTI_RTSR_TR0;
    // Отключить срабатывание по спадающему фронту
    EXTI->FTSR &= ~EXTI_RTSR_TR0;

    // Привязать вывод PA0 к линии прерывания EXTI0
    SYSCFG->EXTICR[0] |= SYSCFG_EXTICR1_EXTI0_PA;

    // Обновить значение системной тактовой частоты
    SystemCoreClockUpdate();
    // Настроить SysTick с периодом прерывания 1 мс
    if (SysTick_Config(SystemCoreClock / 1000))
    {
    }

    // Активировать режим глубокого сна через бит SLEEPDEEP
    SCB->SCR |= SCB_SCR_SLEEPDEEP_Msk;
    // Настроить регистр питания: сбросить бит PDDS
    PWR->CR &= PWR_CR_PDDS;
    // Выбрать режим пониженного потребления standby
    PWR->CR |= PWR_CR_LPDS;

    // Задать схему группировки приоритетов прерываний
    NVIC_SetPriorityGrouping(5);
    // Установить наивысший приоритет для EXTI0
    NVIC_SetPriority(EXTI0_IRQn, 0);
    // Установить приоритет для SysTick (ниже, чем у EXTI0)
    NVIC_SetPriority(SysTick_IRQn, 4);

    // Активировать прерывание от SysTick
    NVIC_EnableIRQ(SysTick_IRQn);
    // Активировать прерывание от внешней линии EXTI0
    NVIC_EnableIRQ(EXTI0_IRQn);

    // Основной цикл: циклическое мигание светодиодов
    while (1)
    {
        DiodPG6();
        DiodPG7();
    }
}
\end{lstlisting}

\newpage

\section{Работа с осциллографом}

Программное измерение периода, частоты сигнала, ширины положительного и отрицательного уровней.

\begin{figure}[H]
\centering
% ВСТАВИТЬ ФОТО 1: программное измерение характеристик сигнала
\includegraphics[width=0.5\textwidth]{Photo/1.jpg}
\caption{Программное измерение характеристик сигнала}
\end{figure}

\begin{figure}[H]
\centering
% ВСТАВИТЬ ФОТО 2: осциллограммы двух каналов с измеренными параметрами
\includegraphics[width=0.5\textwidth]{Photo/2.jpg}
\caption{Осциллограммы двух каналов с измеренными параметрами}
\end{figure}

\newpage

\begin{figure}[H]
\centering
% ВСТАВИТЬ ФОТО 3: попеременное переключение сигналов на двух каналах
\includegraphics[width=0.5\textwidth]{Photo/3.jpg}
\caption{Попеременное переключение сигналов на двух каналах}
\end{figure}

\begin{figure}[H]
\centering
% ВСТАВИТЬ ФОТО 4: спад первого сигнала и нарастание второго
\includegraphics[width=0.5\textwidth]{Photo/4.jpg}
\caption{Спад первого сигнала и нарастание второго}
\end{figure}

\newpage

\section{Результаты работы с осциллографом}

По результатам измерений на осциллографе Tektronix TBS 1052B-EDU получены следующие характеристики сигнала:

\begin{itemize}
    \item Период: 250.6~мс
    \item Частота: 3.990~Гц
    \item Размах: 3.44~В
    \item Амплитуда: 3.28~В
    \item Длительность положительного уровня: 125.2~мс
    \item Длительность отрицательного уровня: 125.3~мс
    \item Скважность: 49.9\%
    \item Время нарастания: 156.0~мкс
    \item Время спада: 152.3~мкс
\end{itemize}

\newpage

% ============================================================
% ЗАКЛЮЧЕНИЕ
% ============================================================
\addcontentsline{toc}{section}{Заключение}
\section*{Заключение}

В ходе выполнения лабораторной работы освоены основные методы формирования временных интервалов с использованием системного таймера SysTick, а также методы перевода микроконтроллера в режим пониженного энергопотребления. Разработана программа на языке C в среде Keil~µVision5, обеспечивающая попеременное включение светодиодов PG6 и PG7 с заданными характеристиками, а также обработку нажатия кнопки «WAKEUP» для перехода МК в спящий режим и его пробуждения. Закреплены навыки работы с осциллографом и оценочной платой MCBSTM32F200 в качестве измерительного генератора.

\newpage

\end{document}